% Complete independent mathematical module. Equation prefix FGM.
% Requires amsmath,amssymb,amsthm,amscd,mathrsfs,hyperref and theorem/proposition.
\section{The ungraded pointed unit and the full supported absolute base}
\label{sec:fgm}

The source used in this section is the pointed multiplicative monoid
\begin{equation}\tag{FGM1}
 M=\{0,\tau_{\mathbf F_1}\},\qquad
 0x=0,\quad \tau_{\mathbf F_1}x=x\quad(x\in M).
\end{equation}
No parity is assigned to $\tau_{\mathbf F_1}$. The parity of an integer
is retained as its residue in $\mathbb Z/2\mathbb Z$. The multiplicative
map below is defined independently of a global addition on a space
containing both the integers and $\tau_{\mathbf F_1}$. We calculate both
the count of occurrences and its idempotent collapse explicitly.

The human sources are Alain Connes and Caterina Consani,
\href{https://arxiv.org/abs/1502.05585v2}{\emph{Absolute algebra and
Segal's $\Gamma$-rings: au dessous de $\overline{\operatorname{Spec}
\mathbb Z}$}}, original \nolinkurl{F1_corr.tex}, equations
\texttt{functsum}, Definition \texttt{defnsalg}, Propositions
\texttt{mono2sss}, \texttt{monomono}, and Lemma \texttt{sssalg}
(lines 360--560 read for this derivation); and their
\href{https://arxiv.org/abs/2606.06604v1}{\emph{On the Absolute Geometry
of Spec Z and the Fargues--Fontaine curve}}, original \nolinkurl{FF.tex},
Section \texttt{sect9.1}, Proposition \texttt{prop12}, and Proposition
\texttt{prop:monoid\_structure\_morphisms} (lines 140--154 and
308--410 read, including the source's exponent-zero Frobenius).
The original coefficient object is the complete lattice-split construction
in \emph{Split Support Geometry}, source byline \emph{The Clankers},
Version 11, Definition \texttt{def:lattice-split} and Theorem
\texttt{thm:lattice-normal-form},
\href{https://github.com/KokunoYumeto/zeta-function-research-reader/blob/a2851f9eb352e7e4376bc629de86fa4ed9c1c434/workbenches/splitzero-tandem/foundations/split-support-geometry-v11/split_support_geometry_arithmetic_curve_v11.tex}{complete pinned proof source}.
All coefficient and receiving-map arguments used here are given below.

\subsection{The original coefficient carrier and its exact base map}
Let $L$ be a bounded distributive lattice with bottom $0_L$ and top $1_L$;
the one-element lattice is allowed. For a nonzero commutative unital
ring $R$ retain the entire carrier
\begin{equation}\tag{FGM2}
 \begin{gathered}
 S_R=G_L(R)=\{z_\lambda=(0,\lambda):\lambda\in L\}
            \cup\{\widehat a=(a,1_L):a\in R\},\\
 (a,\lambda)+(b,\mu)=(a+b,\lambda\vee\mu),\qquad
 (a,\lambda)(b,\mu)=(ab,\lambda\wedge\mu),\\
 \tau_{\rm old}=z_{0_L},\qquad e=z_{1_L}=\widehat0,
 \qquad 1_{S_R}=\widehat1.
 \end{gathered}
\end{equation}
This is a unital semiring. Indeed a nonzero amplitude in a sum forces
top support, and a product with a support below top has zero amplitude.
Thus the carrier is closed; associativity and commutativity follow
coordinatewise, and distributivity follows from that of $R$ and $L$.
The displayed zero and unit follow by direct substitution.

\begin{proposition}[The unique pointed unital map]
The unique pointed unital multiplicative map to the original coefficient
object is
\begin{equation}\tag{FGM3}
 \iota_R:M\longrightarrow(S_R,\cdot),\qquad
 0\longmapsto\tau_{\rm old},\quad
 \tau_{\mathbf F_1}\longmapsto\widehat1.
\end{equation}
It is injective. Its image is the two-element pointed multiplicative
submonoid $\{\tau_{\rm old},\widehat1\}$; every $z_\lambda$ remains
present in its codomain.
\end{proposition}
\begin{proof}
A pointed map must have the displayed value at $0$, and a unital map
must have the displayed value at $\tau_{\mathbf F_1}$. Products involving
$0$ map to the absorbing $\tau_{\rm old}$; the remaining product maps
to $\widehat1\widehat1=\widehat1$. The images are different because
their amplitudes are $0_R$ and $1_R$.
\end{proof}
In particular the new source unit is not sent to the old absorber.
For any multiplicative map $f$ preserving the old multiplication,
the assignment $f(\tau_{\rm old})=1$ forces
$f(x)=f(\tau_{\rm old})f(x)=f(\tau_{\rm old}x)=1$ for every $x$.
When $f$ also preserves zero, its target has $0=1$ and is the
one-element semiring. This proves the exact collapse required by that
assignment, rather than denying a relation between the objects.

The integer $1$ can be retained as distinct from the new global unit
already at the multiplicative level. Define
\begin{equation}\tag{FGM4}
 \mathscr U=\mathbb Z\amalg\{\tau_{\mathbf F_1}\},\qquad
 \tau_{\mathbf F_1}u=u,\quad nm=(nm)_{\mathbb Z}.
\end{equation}
Its zero is the integer zero, and associativity follows by deleting
factors equal to the new unit and multiplying the remaining integers.
Here $1_{\mathbb Z}\tau_{\mathbf F_1}=1_{\mathbb Z}\ne
\tau_{\mathbf F_1}$, so $1_{\mathbb Z}$ is not its global unit.
There is an exact pointed unital multiplicative map
\begin{equation}\tag{FGM5}
 j:\mathscr U\longrightarrow S_{\mathbb Z},\qquad
 j(0)=\tau_{\rm old},\quad j(n)=\widehat n\ (n\ne0),\quad
 j(\tau_{\mathbf F_1})=\widehat1.
\end{equation}
To check multiplication, a zero factor gives the absorber on both
sides; two nonzero integers have nonzero product in $\mathbb Z$;
and a new-unit factor acts identically. Its sole non-singleton fibre
is $\{1_{\mathbb Z},\tau_{\mathbf F_1}\}$. This construction specifies
no mixed addition and assigns no parity to its adjoined unit.

\subsection{Occurrence counts and the two exact additive defects}
Write $\mathbb B=\{0,\tau_{\mathbf F_1}\}$ for the Boolean semiring
whose multiplication is FGM1 and whose additional rule is
$\tau_{\mathbf F_1}+\tau_{\mathbf F_1}=\tau_{\mathbf F_1}$.
This is an explicitly defined receiving quotient, not an assertion
that every expression in the user's larger interaction uses its sum.
There are exact semiring maps
\begin{equation}\tag{FGM6}
 \begin{gathered}
 c:\mathbb N\longrightarrow\mathbb B,
 \qquad c(0)=0,\quad c(n)=\tau_{\mathbf F_1}\ (n>0),\\
 r:\mathbb N\longrightarrow S_{\mathbb Z},
 \qquad r(0)=\tau_{\rm old},\quad r(n)=\widehat n\ (n>0).
 \end{gathered}
\end{equation}
For $c$, a sum is positive precisely when at least one summand is
positive, and a product is positive precisely when both factors are.
For $r$, zero cases use the global absorber, and positive cases use
$\widehat n+\widehat m=\widehat{n+m}$ and
$\widehat n\widehat m=\widehat{nm}$. These checks prove all operations,
including both zeros and units. The map $r$ is injective.
Moreover $\mathbb N$ is the occurrence receiver freely generated by a
multiplicative unit in semirings: any unital semiring map must send
$n$ to the sum of $n$ copies of its target unit, and these assignments
preserve addition and multiplication by expansion of the finite sums.

The full defect of factoring this receiver through $c$ is the pair
relation
\begin{equation}\tag{FGM7}
 c(m)=c(n)\quad(m,n>0),\qquad
 r(m)=\widehat m,\quad r(n)=\widehat n,\qquad
 r(m)=r(n)\ \Longleftrightarrow\ m=n.
\end{equation}
In particular $r$ does not factor through $c$, even as a set map.
The distinct evaluations of one and two occurrences are retained as
\begin{equation}\tag{FGM8}
 \begin{gathered}
 \iota_{\mathbb Z}(\tau_{\mathbf F_1}+_{\mathbb B}
       \tau_{\mathbf F_1})\widehat1=\widehat1,\\
 \iota_{\mathbb Z}(\tau_{\mathbf F_1})\widehat1+
 \iota_{\mathbb Z}(\tau_{\mathbf F_1})\widehat1=\widehat2,
 \qquad r(1+1)=\widehat2.
 \end{gathered}
\end{equation}
Thus the expanded receiver gives $2$ in amplitude. Applying the
Boolean collapse first gives $1$. FGM8 computes the operation defect
of the specified collapse; it does not apply a distributive law in
a source where that law has not been specified.

There is a second, different defect when integer cancellation is
combined with a shared global zero. The supported arithmetic inclusion
\begin{equation}\tag{FGM9}
 k:\mathbb Z\longrightarrow S_{\mathbb Z},\qquad k(n)=\widehat n
\end{equation}
preserves the two binary operations and the unit, with its arithmetic
zero $k(0)=e$. It does not preserve the global zero when $0_L\ne1_L$.
The pointed map $j|_{\mathbb Z}$ of FGM5 instead preserves zero and
multiplication, but has precisely the nonzero cancellation defects
\begin{equation}\tag{FGM10}
 j(n+(-n))=\tau_{\rm old},\qquad
 j(n)+j(-n)=e\quad(n\ne0).
\end{equation}
For all other pairs of integers it preserves addition: a zero
summand uses the absorber, and two nonzero summands of nonzero sum
are covered by FGM2. This exhausts the cases.

\begin{theorem}[The two exact congruences]
For $S=S_{\mathbb Z}$ the smallest semiring congruences imposing,
respectively, the defects in FGM10 and FGM8 are exactly
\begin{equation}\tag{FGM11}
 \begin{aligned}
 S/(e\sim\tau_{\rm old})&\ \simeq\ \mathbb Z,
 &q(a,\lambda)&=a,\\
 S/(\widehat2\sim\widehat1)&\ \simeq\ L,
 &\chi(a,\lambda)&=\lambda.
 \end{aligned}
\end{equation}
The first quotient identifies every $z_\lambda$ and preserves every
integer amplitude. The second preserves every lattice label and
identifies all top-supported integer amplitudes with $e$. Imposing
both relations produces the one-element semiring.
\end{theorem}
\begin{proof}
Both $q$ and $\chi$ preserve zero, unit and both operations by FGM2.
The relation $e\sim\tau_{\rm old}$ implies, for every $\lambda$,
$z_\lambda=e z_\lambda\sim\tau_{\rm old}z_\lambda
=\tau_{\rm old}$. These are all fibres of $q$ with more than one
element; nonzero amplitudes already have unique support $1_L$.
Consequently the generated congruence is exactly equality under $q$.

Add $\widehat{-1}$ to $\widehat2\sim\widehat1$. This gives
$\widehat1\sim e$. Multiplication by $\widehat n$ gives
$\widehat n\sim e$ for every integer $n$. These are all fibres
of $\chi$ with more than one element. Since $\chi$ itself respects
the generating relation, it does not identify two different labels;
thus equality under $\chi$ is exactly the generated congruence.
Finally the two relations together give $1_S\sim e\sim0_S$,
which makes every element equal to zero by multiplication.
\end{proof}
The ungraded pointed base therefore maps into the original full
coefficient object without this quotient. Extending that particular
map to preserve Boolean addition has the exact receiving quotient $L$.
Assigning its shared arithmetic zero to the original global zero has
the different exact receiving quotient $\mathbb Z$.

\subsection{The full spectrum, all stalks, and the sheaf unit}
Let $X=\operatorname{Spec}S$ mean the spectrum of proper additive
semiring prime ideals, and $Y=\operatorname{Spec}\mathbb Z$.
A lattice prime ideal $J\subsetneq L$ is a lower set closed under
finite joins, containing $0_L$, with
$\lambda\wedge\mu\in J$ implying $\lambda\in J$ or $\mu\in J$.
The complete point list is
\begin{equation}\tag{FGM12}
 P_J=\{z_\lambda:\lambda\in J\},\qquad
 Q_{\mathfrak p}=q^{-1}(\mathfrak p),\qquad
 \mathfrak p=(0)\ \hbox{or}\ (p),\qquad P_J\subset Q_{(0)}.
\end{equation}
To prove completeness, an ideal containing a nonzero $\widehat n$
also contains $\widehat{-n}$, hence $e$ and all $z_\lambda=e z_\lambda$.
An ideal containing $e$ is the full inverse image of an ideal of
$\mathbb Z$: its top-supported amplitudes are closed under sums,
negatives, and multiplication by integers. Primality is exactly
primality of that ring ideal. An ideal not containing $e$ consists
only of zero labels; addition is join, multiplication by $z_\mu$
is meet, and the ideal and prime axioms are exactly those of a
proper lattice prime $J$. Conversely these formulas verify all
ideal and prime axioms for the displayed list. The case of all zero
labels is $Q_{(0)}$, not an omitted additional support point.

Put $A=\{P_J\}$. Direct membership in the basic opens gives
\begin{equation}\tag{FGM13}
 \begin{gathered}
 D_S(z_\lambda)=\{P_J:\lambda\notin J\},\qquad D_S(e)=A,\\
 D_S(\widehat n)=A\ \cup\ \{Q_{\mathfrak p}:n\notin\mathfrak p\}
 \quad(n\ne0).
 \end{gathered}
\end{equation}
The continuous map $r_X:X\to Y$ sends $P_J,Q_{(0)}$ to the generic
point $\eta$ and $Q_{(p)}$ to $p$. Indeed for every nonzero integer
$n$, the inverse image of $D_{\mathbb Z}(n)$ is $D_S(\widehat n)$;
these opens and the empty set form a basis. For $n=0$ its inverse
image is empty, whereas $D_S(\widehat0)=A$: this exceptional value
is retained. The map $i:Y\to X$, $\mathfrak p\mapsto Q_{\mathfrak p}$,
identifies $Y$ with $V(e)$, and $r_Xi=\operatorname{id}_Y$.

Let $\mathcal O_X$ be the sheaf of locally represented fractions of
$S$. For each lattice prime $J$ put
\begin{equation}\tag{FGM14}
 L_J=L[(L\setminus J)^{-1}],\qquad
 [\lambda]=[\mu]\ \Longleftrightarrow\
 \exists\nu\notin J:\lambda\wedge\nu=\mu\wedge\nu.
\end{equation}
Its operations are induced join and meet. The exact stalks and
generizations are
\begin{equation}\tag{FGM15}
 \begin{gathered}
 (\mathcal O_X)_{Q_{(p)}}=G_L(\mathbb Z_{(p)}),\quad
 (\mathcal O_X)_{Q_{(0)}}=G_L(\mathbb Q),\quad
 (\mathcal O_X)_{P_J}=L_J,\\
 G_L(\mathbb Z_{(p)})\longrightarrow G_L(\mathbb Q),\quad
 (a,\lambda)\longmapsto(a,\lambda),\qquad
 G_L(\mathbb Q)\longrightarrow L_J,\quad(a,\lambda)\longmapsto[\lambda].
 \end{gathered}
\end{equation}
At an arithmetic point the denominators are exactly
$\widehat n$ with $n\notin\mathfrak p$. Fractions map to
$(a/n,\lambda)$; their equality is the ordinary ring localization
equality together with equality of the unchanged support labels.
This proves the first two stalk formulas. At $P_J$ the denominator
$e$ is allowed. Since $e^2=e$, inverting it gives $e=1$ and
$(a,\lambda)=e(a,\lambda)=z_\lambda$. Inverting the remaining
$z_\nu$, $\nu\notin J$, gives FGM14 and the third formula.
An invertible idempotent is $1$, so fractions in $L_J$ are represented
by their numerators, with exactly the equality relation FGM14.
These constructions prove all displayed maps. For $J\subseteq K$
the further map $L_K\to L_J$ is $[\lambda]_K\mapsto[\lambda]_J$.
The bottom and top in $L_J$ differ: their equality would require
$\nu\notin J$ with $\nu=0_L$, contrary to $0_L\in J$.
Inverting $e$ therefore identifies it with the unit, not with zero.

Retain the Connes--Consani source sheaf with its original stalks:
\begin{equation}\tag{FGM16}
 \begin{gathered}
 \mathcal F=\Theta^{-1}\mathcal O_{\mathbf F_1},\qquad
 \mathcal O_{\mathbf F_1}=\mathbf F_1[T],\qquad
 \Theta(\eta)=[\{0\}],\quad\Theta(p)=[\mathbb Z[1/p]],\\
 \mathcal F_\eta=\mathbf F_1,\qquad
 \mathcal F_p=\mathbf F_1[T^{H_p^+}],\qquad
 H_p^+=\mathbb Z[1/p]\cap[0,\infty),\qquad
 \mathcal F_L=r_X^{-1}\mathcal F.
 \end{gathered}
\end{equation}
The inverse image here is the sheafification of the presheaf colimit
over opens $V\subseteq Y$ with $W\subseteq r_X^{-1}V$.
Taking a stalk gives the colimit over neighborhoods of $r_X(x)$,
so $(\mathcal F_L)_x=\mathcal F_{r_X(x)}$. Thus every $P_J$ and
$Q_{(0)}$ has stalk $\mathbf F_1$, and every $Q_{(p)}$ has the
displayed monoid algebra. The source includes the endomorphism
$\operatorname{Fr}_0(T^h)=T^0=1$, fixing its separate absorbing zero.
Every generization commutes with it; on $\mathbf F_1$ it is the
identity. Therefore the exact finite-to-generic map is
\begin{equation}\tag{FGM17}
 \epsilon_p:\mathcal F_p\longrightarrow\mathbf F_1,
 \qquad \epsilon_p(0)=0,\qquad\epsilon_p(T^h)=1\quad(h\ge0).
\end{equation}
Indeed $g(T^h)=\operatorname{Fr}_0g(T^h)=g(T^0)=1$.
All maps between the generic $\mathbf F_1$ stalks are the identity.

\begin{theorem}[The full-support spherical map]
There is exactly one unital spherical-algebra sheaf map
\begin{equation}\tag{FGM18}
 \beta:\mathcal F_L\longrightarrow H\mathcal O_X.
\end{equation}
Its degree-one values are
\begin{equation}\tag{FGM19}
 \begin{array}{c|cc}
 \text{stalk}&\text{source absorbing zero}&\text{source unit or }T^h\\ \hline
 Q_{(p)}&\tau_{\rm old}&\widehat1\quad(h\in H_p^+)\\
 Q_{(0)}&\tau_{\rm old}&\widehat1\\
 P_J&[0_L]&[1_L].
 \end{array}
\end{equation}
Every support label remains in the target stalk or its exact
localization class from FGM14.
\end{theorem}
\begin{proof}
For a semiring $B$, $(HB)(k_+)=B^k$; a pointed map acts by sums over
non-basepoint fibres, with empty sum $0_B$, and multiplication takes
$(b_i),(c_j)$ to $(b_ic_j)_{i,j}$. For a pointed monoid $N$,
the spherical monoid algebra in degree $k_+$ is $N\wedge k_+$.
A pointed unital monoid map $f:N\to B$ induces
\begin{equation}\tag{FGM20}
 (n,j)\longmapsto(0_B,\ldots,0_B,f(n),0_B,\ldots,0_B),
\end{equation}
with the indicated value in coordinate $j$ and the basepoint sent
to the all-zero tuple. A pointed map moves the single occupied
coordinate or sends it to the basepoint; the fibre-sum rule does the
same. Products agree by multiplicativity of $f$. Conversely the maps
$1_+\to k_+$ selecting a coordinate force FGM20 from its degree-one
component. This proves the needed source adjunction explicitly.

The original $\operatorname{Fr}_0$ factors through the constant
spherical algebra $\mathbf F_1$ and commutes with all source
Frobenius maps. Pulling back this factorization gives the sheaf
augmentation $\mathcal F_L\to\underline{\mathbf F_1}$. Compose
with the spherical unit map of $H\mathcal O_X$, given by FGM20
and the stalk units of FGM15. This constructs $\beta$ and proves
FGM19 globally. For uniqueness its component at $Q_{(0)}$ is the
unique unit map. The arithmetic generization in FGM15 is injective;
compatibility with FGM17 forces every $T^h$ at $Q_{(p)}$ to have
image $\widehat1$. At $P_J$ the source is $\mathbf F_1$, whose
unital map is unique. Formula FGM20 determines all degrees; equal
maps on every stalk are equal sheaf maps.
\end{proof}
There is no asserted Boolean sum in the spherical source unit.
In degree $2_+$ the sphere is $\{0,1,2\}$. Its two coordinate
projections give $(0,0)$, $(1,0)$, and $(0,1)$ respectively;
none gives $(1,1)$. Thus there is no degree-two source element
encoding two occupied unit coordinates. In $H\mathbb B$ such
an element does exist and folds to $\tau_{\mathbf F_1}$; in
$HS_{\mathbb Z}$ it folds to $\widehat2$. FGM20 and FGM8
give the exact relation between these constructions.

\subsection{Every zero branch and the exact parity reduction}
For completeness retain all algebraic finite-prime points with
target $HS_{\mathbb C}$. A character is a multiplicative map
$\xi:H_p^+\to S_{\mathbb C}$ with $\xi(0)=\widehat1$;
the extra absorbing zero of the monoid algebra maps separately
to $\tau_{\rm old}$. The entire set is
\begin{equation}\tag{FGM21}
 \begin{gathered}
 \mathcal U_p=\{(u_n)_{n\ge0}\in(\mathbb C^\times)^{\mathbb N}:
                         u_{n+1}^{p}=u_n\},\qquad
 \operatorname{Hom}_{\mathbf F_1}(\mathcal F_p,HS_{\mathbb C})
       =\mathcal U_p\ \amalg\ \{Z_\lambda:\lambda\in L\},\\
 \xi_u(a/p^n)=\widehat{u_n^a}\quad(a\ge0),\qquad
 Z_\lambda(0)=\widehat1,\quad Z_\lambda(h)=z_\lambda\quad(h>0).
 \end{gathered}
\end{equation}
To prove completeness set $b_n=\xi(p^{-n})$, so $b_{n+1}^p=b_n$.
Nonzero complex amplitudes stay nonzero under positive powers.
Thus either every $b_n$ has nonzero amplitude, giving exactly the
displayed sequence $u$, or every $b_n$ is a zero label. In the
second case $z_\lambda^p=z_\lambda$, hence the labels are all one
fixed $\lambda$. Every positive exponent is $a/p^n$ with $a>0$,
so these cases determine the full character. Conversely the
formulas respect a change of denominator, exponent addition,
the unit and the separate zero, and therefore define all points.
Their multiplication retains every label:
\begin{equation}\tag{FGM22}
 \xi_u\xi_v=\xi_{uv},\qquad \xi_uZ_\lambda=Z_\lambda,
 \qquad Z_\lambda Z_\mu=Z_{\lambda\wedge\mu}.
\end{equation}
The unit point is the constant unit sequence. The absorber point
is $Z_{0_L}$, with value $\widehat1$ at exponent zero, not an
everywhere-zero function. The amplitude map identifies all
$Z_\lambda$ with the single zero branch and preserves each $u$.
The generic point map induced by FGM17 selects the constant unit
sequence, not a zero branch. For every integer $m\ge1$ source
Frobenius gives $u\mapsto(u_n^m)$ and fixes each $Z_\lambda$;
for $m=0$ it sends every point to the unit point, with the extra
source absorbing zero still sent to $\tau_{\rm old}$.
These assertions follow by evaluating at $mh$, including $h=0$.

\begin{proposition}[Parity with full support]
There is a surjective unital semiring map
\begin{equation}\tag{FGM23}
 \pi_2:G_L(\mathbb Z)\longrightarrow G_L(\mathbb F_2),
 \qquad (n,\lambda)\longmapsto(\overline n,\lambda).
\end{equation}
Its exact fibres are
\begin{equation}\tag{FGM24}
 \begin{aligned}
 \pi_2^{-1}(\widehat1)&=\{\widehat n:n\text{ odd}\},\\
 \pi_2^{-1}(e)&=\{\widehat n:n\text{ even}\},\\
 \pi_2^{-1}(z_\lambda)&=\{z_\lambda\}\quad(\lambda\ne1_L).
 \end{aligned}
\end{equation}
Consequently the zero-amplitude fibre contains every zero label and
every even integer; the nonzero-amplitude fibre consists exactly of
the odd integers. All reduced nonzero amplitudes are odd. There is
no assigned parity for $\tau_{\mathbf F_1}$ in these statements.
\end{proposition}
\begin{proof}
Reduction respects ring operations, while both lattice coordinates
in FGM2 are unchanged. It therefore respects the full operations,
zeros and units; each target element has an indicated preimage.
All nonzero source amplitudes have top support, and the only
nonzero residue in $\mathbb F_2$ is $1$. These observations prove
each fibre formula. On the supported integer copy $k$, the
amplitude composite is exactly the intrinsic residue
$n\mapsto\overline n$; its being a numeric invariant does not
assign a parity to the different source unit in FGM1.
\end{proof}

At the arithmetic stalk over $2$ the reduction is
\begin{equation}\tag{FGM25}
 G_L(\mathbb Z_{(2)})\longrightarrow G_L(\mathbb F_2),\qquad
 (a/b,\lambda)\longmapsto(\overline a\,\overline b^{-1},\lambda),
 \qquad b\text{ odd}.
\end{equation}
Cross multiplication proves it is well-defined, since every
allowed denominator reduces to $1$; the operation checks are
those of FGM23. For $p\ne2$ there is no unital semiring map from
$G_L(\mathbb Z_{(p)})$ to $G_L(\mathbb F_2)$ extending FGM23,
and none from $G_L(\mathbb Q)$: $\widehat2$ is invertible in
either source, while its required image $e$ is not invertible.
Indeed every product with $e$ has zero amplitude, so it cannot
be $\widehat1$ in the target. This proves the exact local domain
of parity reduction.

At each pure-support stalk the precise compatibility instead is
the commuting square
\begin{equation}\tag{FGM26}
 \begin{CD}
 G_L(\mathbb Z) @>{\pi_2}>> G_L(\mathbb F_2)\\
 @V{(n,\lambda)\mapsto[\lambda]}VV
 @VV{(a,\lambda)\mapsto[\lambda]}V\\
 L_J @= L_J .
 \end{CD}
\end{equation}
Both routes preserve the entire resulting label class and discard
amplitude by the proved map FGM15. They identify all integer
images with $[1_L]$; the label equality remains exactly FGM14.
Thus parity is not recoverable from that support quotient.
In particular FGM8 reduces over $2$ to the distinct pair
$(\widehat1,e)$, whereas its two entries both become $[1_L]$
in a pure-support stalk. The image of the new base unit is always
the appropriate target unit in FGM19; its ungraded source
designation remains unchanged.

